\subsection{The dual reduction for a hierarchy}
\label{app:hierarchy-dual}

This section proves the matrix-to-embedding reduction needed for the
hierarchy estimate.  The marked nodes may be any subset of the nonroot
nodes.  In fact, the only hierarchical property used in this reduction is
that an edge belongs to at most two of the sets $O(t)$.

Let $G=(V,E)$ be a connected loopless multigraph, with $n=|V|\geq2$
and $m=|E|$.  Fix an orientation of every edge, write $b_e$ for its
incidence column, and put
\[
 B=[b_e]_{e\in E},\qquad L_G=BB^{\mathsf T},\qquad
 \mathcal C=\{S:\varnothing\ne S\subsetneq V\},\qquad
 d_S=|\delta_G(S)|.
\]
Thus $d_S=\mathbf1_S^{\mathsf T}L_G\mathbf1_S>0$ for $S\in\mathcal C$.
Let $\mathcal M$ be the set of marked nodes, and assume that
$O(t)\ne\varnothing$ for $t\in\mathcal M$.  Define the value of
the shortcut program by an infimum:
\begin{equation}
\label{hd:primal}
 p_*=
 \inf_{\substack{D\succ0\\
       \mathbf1_S^{\mathsf T}D\mathbf1_S\leq d_S\ (S\in\mathcal C)}}
 \ \max_{t\in\mathcal M}
 \frac1{|O(t)|}\sum_{e\in O(t)} b_e^{\mathsf T}D^{-1}b_e.
\end{equation}
Here and throughout this section, $D\succ0$ means positive definite on
the full space $\mathbb R^V$.  No constraint for $S=V$ is imposed.
When $\mathcal M=\varnothing$, set the objective and $p_*$ equal to zero.

For a binary matrix $X\in\{0,1\}^{h\times V}$, let
\[
 \mathcal E(X)=\sum_{e\in E}\|Xb_e\|_2^2.
\]
A matrix $U\in\mathbb R^{E\times h}$ will be called
\emph{row semiorthogonal} if $UU^{\mathsf T}=I_m$; its row indexed by
$e$ is denoted by $u_e^{\mathsf T}$.  Thus the vectors $u_e\in\mathbb R^h$
are orthonormal, and this convention requires $h\geq m$.  The dimension
$h$ is otherwise arbitrary.  Zero rows can always be added to $X$.

\Needspace*{3\baselineskip}
\begin{lemma}[Dual reduction]
\label{lem:hierarchy-dual}
With the preceding notation, set
\begin{equation}
\label{hd:ratio}
 Q_*=
 \sup_{\substack{h\geq m,\ X\in\{0,1\}^{h\times V},\ \mathcal E(X)>0\\
                  U\in\mathbb R^{E\times h},\ UU^{\mathsf T}=I_m}}
 \frac{\displaystyle
       \sum_{t\in\mathcal M}\frac1{|O(t)|}
       \left(\sum_{e\in O(t)}\langle u_e,Xb_e\rangle\right)^2}
      {\mathcal E(X)}.
\end{equation}
Then
\begin{equation}
\label{hd:comparison}
 p_*\leq Q_*\leq2p_*.
\end{equation}
Consequently, if every ratio in \eqref{hd:ratio} is at most $K$, then for
every $\eta>0$ there is a full positive definite matrix $D$ satisfying
all the cut constraints in \eqref{hd:primal} and all the marked-node
average-resistance bounds with right side $K+\eta$.
\end{lemma}

\begin{proof}
The assertion for an empty marked set follows by taking a sufficiently
small positive multiple of the identity, so assume $\mathcal M\ne\varnothing$.
We give the finite-dimensional duality argument as well as the matrix
and embedding steps.

\paragraph{Separation in finite dimension.}
We use the following elementary geometric fact: if $\mathcal U$ is a
nonempty open convex subset of a Euclidean space and $a\notin\mathcal U$,
there is a nonzero vector $v$ such that
$\langle v,x-a\rangle\geq0$ for every $x\in\mathcal U$.
Here is a proof.  Choose $y\in\mathcal U$ and put
$a_j=a+(a-y)/j$.  None of the points $a_j$ belongs to
$K=\overline{\mathcal U}$.  Otherwise $a$, being a strict convex
combination of $y\in\mathcal U$ and $a_j\in K$, would belong to
$\mathcal U$: an open ball about $y$, convexly combined with points of
$\mathcal U$ approaching $a_j$, gives an open ball about $a$ in
$\mathcal U$.  Choose a nearest point $z_j\in K$ to $a_j$.
Such a point exists by minimizing distance on a sufficiently large
closed ball.  Convexity and differentiation of squared distance along
the segment from $z_j$ to any $x\in K$ give
\[
 \langle z_j-a_j,x-z_j\rangle\geq0.
\]
Hence, for $v_j=(z_j-a_j)/\|z_j-a_j\|_2$,
$\langle v_j,x-a_j\rangle\geq\|z_j-a_j\|_2\geq0$.
A convergent subsequence of the unit vectors $v_j$ has a unit limit
$v$, which has the stated property since $a_j\to a$.

\paragraph{Strong duality for this program.}
Put
\[
 A_t=\frac1{|O(t)|}\sum_{e\in O(t)}b_eb_e^{\mathsf T},\qquad
 f_t(D)=\operatorname{tr}(A_tD^{-1}),\qquad
 g_S(D)=\mathbf1_S^{\mathsf T}D\mathbf1_S-d_S.
\]
Each $f_t$ is convex on $D\succ0$.  Indeed, completing the square gives
\[
 b^{\mathsf T}D^{-1}b
   =\sup_{x\in\mathbb R^V}
      \bigl(2b^{\mathsf T}x-x^{\mathsf T}Dx\bigr),
\]
a supremum of affine functions of $D$.  The functions $g_S$ are affine.
Choose
\[
 0<\varepsilon_0<\min_{S\in\mathcal C}\frac{d_S}{|S|},
 \qquad D_0=\varepsilon_0 I.
\]
Then all cut constraints are strict at $D_0$.  In particular
$0\leq p_*<\infty$.

Consider the open convex set of vectors $(a,b,s)$ for which there are
$D\succ0$ and $r\in\mathbb R$ such that
\[
 f_t(D)-r<a_t\quad(t\in\mathcal M),\qquad
 g_S(D)<b_S\quad(S\in\mathcal C),\qquad r<s.
\]
The point $(0,0,p_*)$ is outside this set, since membership would give a
feasible $D$ with every $f_t(D)<p_*$.  The separation fact just proved
therefore supplies nonzero coefficients $(\lambda,y,\alpha)$ such that
\[
 \lambda\cdot a+y\cdot b+\alpha s\geq\alpha p_*
\]
throughout the set.  Each coefficient is nonnegative, because any one
coordinate of a point in the set can be increased arbitrarily.
Moreover $\alpha>0$.  To see this, choose
$r_0>\max_t f_t(D_0)$; there is then a point of the set with all the
$a_t$ and $b_S$ negative.  If $\alpha=0$, this would contradict the
separation inequality and the nonzero, nonnegative vector $(\lambda,y)$.
Normalize $\alpha=1$ and take limits down to the displayed strict
inequalities.  For every $D\succ0$ and every $r\in\mathbb R$ this gives
\[
 \sum_t\lambda_tf_t(D)+\sum_Sy_Sg_S(D)
       +\left(1-\sum_t\lambda_t\right)r\geq p_*.
\]
Since $r$ is unrestricted, $\sum_t\lambda_t=1$.

For any nonnegative $\lambda$ with this sum and any $y\geq0$, weak
duality gives the reverse bound: at a feasible $D$, the weighted
objective plus the weighted constraint residuals is at most
$\max_t f_t(D)$.  Consequently, writing
\[
 A(\lambda)=\sum_t\lambda_tA_t,\qquad
 Z(y)=\sum_{S\in\mathcal C}y_S\mathbf1_S\mathbf1_S^{\mathsf T},\qquad
 d(y)=\sum_{S\in\mathcal C}y_Sd_S,
\]
we have proved
\begin{equation}
\label{hd:lagrange}
 p_*=
 \sup_{\substack{\lambda\geq0,\ \sum_t\lambda_t=1\\y\geq0}}
 \inf_{D\succ0}
 \left\{\operatorname{tr}(A(\lambda)D^{-1})
       +\operatorname{tr}(Z(y)D)-d(y)\right\}.
\end{equation}
This argument neither assumes attainment of the primal infimum nor
requires a uniform positive lower bound on the eigenvalues of its
feasible matrices.

\paragraph{The matrix infimum, including singular matrices.}
For any symmetric positive semidefinite $A,Z$, we claim that
\begin{equation}
\label{hd:matrix-inf}
 \inf_{D\succ0}
 \bigl\{\operatorname{tr}(AD^{-1})+\operatorname{tr}(ZD)\bigr\}
 =2\operatorname{tr}\bigl((Z^{1/2}AZ^{1/2})^{1/2}\bigr).
\end{equation}
First suppose $A,Z\succ0$, and set
$H=Z^{1/2}DZ^{1/2}$ and $C=Z^{1/2}AZ^{1/2}$.  The difference between
the left-hand objective and $2\operatorname{tr}(C^{1/2})$ equals
\[
 \operatorname{tr}(CH^{-1})+\operatorname{tr}(H)
       -2\operatorname{tr}(C^{1/2})
   =\|H^{-1/2}C^{1/2}-H^{1/2}\|_{\mathrm F}^2\geq0.
\]
It is zero when $H=C^{1/2}$, proving the identity in this case.
For general $A,Z\succeq0$, replace them temporarily by
$A_\varepsilon=A+\varepsilon I$ and
$Z_\varepsilon=Z+\varepsilon I$.  The positive definite identity,
applied at any fixed $D$, and then continuity of positive semidefinite
square roots give the lower bound in \eqref{hd:matrix-inf} as
$\varepsilon\downarrow0$.  Conversely, the unperturbed objective is
at most the perturbed objective at every $D$, so its infimum is at
most the known perturbed infimum.  Letting $\varepsilon\downarrow0$
gives the upper bound.  Continuity here also follows directly from the
spectral theorem: along a convergent sequence of positive semidefinite
matrices their square roots are bounded, and every convergent
subsequence of these square roots has the unique positive
semidefinite square root of the limit as its limit.  These
perturbations prove an auxiliary matrix identity; they do not change
the allowed dual multipliers.

For fixed $\lambda,y$, abbreviate the trace of the square root in
\eqref{hd:matrix-inf} by $N(\lambda,y)$.  If $y\ne0$, then $d(y)>0$.
Scaling $y$ by $s\geq0$ changes the dual value to
$2\sqrt{s}\,N(\lambda,y)-s\,d(y)$, whose supremum over $s\geq0$ is
$N(\lambda,y)^2/d(y)$.  The case $y=0$ contributes zero.  Thus
\eqref{hd:lagrange} becomes
\begin{equation}
\label{hd:homogeneous}
 p_*=
 \sup_{\substack{\lambda\geq0,\ \sum_t\lambda_t=1\\y\geq0,\ y\ne0}}
 \frac{N(\lambda,y)^2}{d(y)}.
\end{equation}

\paragraph{From weighted cuts to binary coordinates.}
Let $W=W(\lambda)$ be the diagonal $E\times E$ matrix with entries
\begin{equation}
\label{hd:weights}
 w_e=\left(\sum_{t\in\mathcal M:\ e\in O(t)}
                      \frac{\lambda_t}{|O(t)|}\right)^{1/2}.
\end{equation}
Then $A(\lambda)=BW^2B^{\mathsf T}$.  If a matrix $Y$ has rows
$\sqrt{y_S}\,\mathbf1_S^{\mathsf T}$, then $Y^{\mathsf T}Y=Z(y)$.
The matrices $Z(y)^{1/2}A(\lambda)Z(y)^{1/2}$ and
$WB^{\mathsf T}Z(y)BW$ have the same nonzero eigenvalues: they are
the two products $TT^{\mathsf T}$ and $T^{\mathsf T}T$ with
$T=Z(y)^{1/2}BW$.  It follows that
\[
 N(\lambda,y)=\|YBW\|_*,\qquad
 d(y)=\sum_{e\in E}\|Yb_e\|_2^2,
\]
where $\|M\|_*$ denotes the sum of the singular values of $M$.

For fixed $\lambda$, the ratio in \eqref{hd:homogeneous} is continuous
in $y\geq0$, $y\ne0$.  Approximate $y$ by nonzero nonnegative rational
vectors.  For rational $y$, choose a positive integer $q$ with every
$qy_S$ integral and form a binary matrix $X$ by repeating the row
$\mathbf1_S^{\mathsf T}$ exactly $qy_S$ times.  Then
$X^{\mathsf T}X=qZ(y)$, and both the squared numerator and the
denominator in the ratio acquire the same factor $q$.  Conversely,
any binary $X$ gives integral cut multiplicities after its constant
rows are discarded.  Constant rows contribute zero to $XB$, so they
have no effect on the ratio.  Adding zero rows permits $h\geq m$.
These observations prove the exact identity
\begin{equation}
\label{hd:binary-nuclear}
 p_*=
 \sup_{\substack{h\geq m,\ X\in\{0,1\}^{h\times V},\ \mathcal E(X)>0\\
                  \lambda\geq0,\ \sum_t\lambda_t=1}}
 \frac{\|XBW(\lambda)\|_*^2}{\mathcal E(X)}.
\end{equation}
The limiting argument takes a supremum over all finite dimensions;
it asserts no exact finite binary representation for irrational
cut weights.

\paragraph{Orthonormal edge vectors and the factor two.}
For any $h\times m$ matrix $M$ with $h\geq m$, singular-value
decomposition gives
\begin{equation}
\label{hd:polar}
 \|M\|_*=\max_{UU^{\mathsf T}=I_m}\operatorname{tr}(UM).
\end{equation}
Indeed, write $M=P\Sigma V^{\mathsf T}$ with
$P^{\mathsf T}P=I_m$, $V$ orthogonal, and $\Sigma$ diagonal
nonnegative; zero singular values can be included by completing the
columns of $P$.  For any admissible $U$, each diagonal entry of
$V^{\mathsf T}UP$ is at most one, giving the upper bound
$\operatorname{tr}(UM)\leq\operatorname{tr}(\Sigma)$.
The choice $U=VP^{\mathsf T}$ attains it.

Write $a_e=\langle u_e,Xb_e\rangle$.  Individual rows of $U$ can be
multiplied by $-1$ while preserving $UU^{\mathsf T}=I_m$.
Consequently, when taking either the supremum in \eqref{hd:ratio} or
the supremum of the squared trace in \eqref{hd:polar}, we may restrict
to $a_e\geq0$ for all $e$: replacing $a_e$ by $|a_e|$ cannot decrease
either expression.  Equation \eqref{hd:binary-nuclear} therefore says
\begin{equation}
\label{hd:weighted-rows}
 p_*=
 \sup_{X,U,\lambda}
 \frac{\left(\sum_e w_ea_e\right)^2}{\mathcal E(X)},
 \qquad a_e\geq0,
\end{equation}
with the other conditions on $X,U,\lambda$ as above.

Every edge belongs to at most two parent boundaries $O(t)$: if its
endpoints have lowest common ancestor $s$, only the two children of
$s$ containing its respective endpoints can have that edge in their
parent boundary.  Restricting to marked nodes preserves this property.
Hence, on putting
\[
 v_e=\sum_{t\in\mathcal M:\ e\in O(t)}
                    \sqrt{\lambda_t/|O(t)|},
\]
the elementary inequalities between the Euclidean norm and the sum
of at most two nonnegative numbers give
$w_e\leq v_e\leq\sqrt2\,w_e$.  Because $a_e\geq0$,
\begin{equation}
\label{hd:overlap}
 \left(\sum_e w_ea_e\right)^2
 \leq\left(\sum_e v_ea_e\right)^2
 \leq2\left(\sum_e w_ea_e\right)^2.
\end{equation}
For fixed $X,U$, let
\[
 c_t=\frac1{\sqrt{|O(t)|}}\sum_{e\in O(t)}a_e\geq0.
\]
Then $\sum_e v_ea_e=\sum_t\sqrt{\lambda_t}\,c_t$.
Cauchy--Schwarz, with equality at
$\lambda_t=c_t^2/\sum_s c_s^2$ when the denominator is nonzero,
shows that
\[
 \max_{\lambda\geq0,\ \sum_t\lambda_t=1}
       \left(\sum_e v_ea_e\right)^2
   =\sum_t c_t^2.
\]
The identity is also valid when all $c_t$ vanish.  Taking suprema in
\eqref{hd:overlap} and using \eqref{hd:weighted-rows} proves
$p_*\leq Q_*\leq2p_*$.  Finally, if $Q_*\leq K$, the defining
property of the finite infimum $p_*$ supplies a feasible $D\succ0$
whose maximum marked average is less than $p_*+\eta\leq K+\eta$.
\end{proof}

\paragraph{Dimension and orientation conventions.}
The ratio does not depend on the chosen edge orientations, since a
reversed $b_e$ can be accompanied by a reversed $u_e$.  Our use of
arbitrarily large $h$ also covers the other usual rectangular
semiorthogonal convention.  If $h<m$ and an $m\times h$ matrix has
orthonormal columns, complete its columns to an $m\times m$ orthogonal
matrix and append $m-h$ zero rows to $X$.  Every edge inner product
and the denominator are preserved.  More generally, any matrix with
$UU^{\mathsf T}\preceq I_m$ can be enlarged to the row semiorthogonal
matrix
\[
 \bigl[\,U\quad (I_m-UU^{\mathsf T})^{1/2}\,\bigr]
\]
by appending $m$ zero rows to $X$, again preserving the ratio.
